\section{Desarrollo}
% Metodologia / Desarrollo
Partimos de la obtencion de nuestros datos que nos servian para codificar y generar nuestros individuos por lo cual se hizo una base de datos investigando el plan de estudios
de cada curso considerado para este trabajo los cuales se mencionaron en la seccion anterior.~\ref{cursos}

La resolución del UCTP mediante Algoritmos Genéticos requiere definir una representación clara del problema, así como los operadores evolutivos y una función de aptitud 
capaz de satisfacer las restricciones duras y optimizar las restricciones blandas.

\subsection{Representación del Cromosoma}

Cada individuo de la población representa un horario completo. Para ello, se construyó un cromosoma donde cada gen corresponde a un evento académico (materia--grupo) 
definido en la base de datos. Cada gen se codifica como un par ordenado:

\[
\text{gen}_i = (\text{room\_index}_i, \text{start\_slot}_i)
\]

donde:
\begin{itemize}
    \item \textbf{room\_index} identifica el salón asignado, indexado en la tabla de salones disponibles.
    \item \textbf{start\_slot} representa el bloque de tiempo de inicio en la grilla semanal.
\end{itemize}

Esta representación permite manejar clases de diferente duración, tipos de aula (laboratorio y salón de clases), y múltiples cursos simultáneamente. 
Además, facilita la verificación de cruces, compatibilidad y duración sin necesidad de estructuras intermedias complejas.

\subsection{Datos y Restricciones Incorporadas}

La base de datos utilizada contiene:
\begin{itemize}
    \item Lista de cursos, su duración en bloques y el número de sesiones por semana.
    \item Relación curso--grupo para evitar combinar grupos distintos en un mismo curso.
    \item Tipos de salón requeridos (normal, laboratorio).
    \item Tabla de salones disponibles con su \textit{room\_type} correspondiente.
    \item Número total de bloques de tiempo por día y por semana.
\end{itemize}

A partir de estas tablas se construyeron estructuras auxiliares como:
\begin{itemize}
    \item \texttt{compatible\_rooms[i]}: lista de salones posibles para el evento $i$.
    \item \texttt{duration\_slots[i]}: número de bloques que ocupa cada clase.
\end{itemize}

\subsection{Generación de Individuos}

Para crear la población inicial, cada evento se asigna de la siguiente manera:

\begin{enumerate}
    \item Selección aleatoria de un salón dentro de \texttt{compatible\_rooms}.
    \item Selección aleatoria del \texttt{start\_slot}, asegurando que la clase completa quepa en el día.
    \item Verificación básica para evitar cruces evidentes dentro del mismo individuo.
\end{enumerate}

Este proceso genera soluciones diversas que cumplen mínimamente con las restricciones duras estructurales.

\subsection{Función de Aptitud}

La función de aptitud se define como la suma de penalizaciones por incumplimiento de restricciones. Esto permite mantener una formulación flexible y ajustable. 
Las penalizaciones actuales incluyen:

\begin{itemize}
    \item \textbf{Cruce de dos clases en el mismo salón y mismo bloque} (penalización alta).
    \item \textbf{Uso de un salón incompatible} con el tipo requerido.
    \item \textbf{Clases que exceden el día} por mala asignación de \textit{start\_slot}.
    \item \textbf{Más de dos horas consecutivas}: si un grupo tiene tres o más clases seguidas de la misma materia, se aplica penalización creciente.
    \item \textbf{Distribución semanal deficiente}: cuando un curso concentra todas sus horas en un mismo día.
\end{itemize}

La aptitud total se define como:

\[
\text{fitness}(ind) = \sum_j \text{penalización}_j(ind)
\]

y el objetivo del algoritmo es minimizarla.

\subsection{Operadores Genéticos}

Se emplean los operadores clásicos adaptados a la estructura del horario:

\begin{itemize}
    \item \textbf{Selección por torneo}: favorece individuos con menor penalización.
    \item \textbf{Cruce uniforme por gen}: cada gen del hijo proviene aleatoriamente de uno de los padres, manteniendo la pareja (salón, inicio).
    \item \textbf{Mutación controlada}: se permite cambiar el salón o el bloque de inicio respetando las compatibilidades.
\end{itemize}

Estos operadores permiten explorar el espacio de soluciones sin destruir completamente buenas configuraciones parciales.

\subsection{Construcción del Horario Final}

Una vez que el AG converge a una solución con baja penalización, se utiliza una función de decodificación para transformar el cromosoma en un horario estructurado por día. 
Para cada día, se genera automáticamente:
\begin{itemize}
    \item Una matriz con bloques de tiempo en el eje vertical.
    \item Los salones disponibles en el eje horizontal.
    \item Las clases asignadas en cada celda con su nombre y grupo.
\end{itemize}

Finalmente, se exporta a un archivo Excel donde cada hoja representa un día de la semana. Este archivo también incorpora formatos de color para distinguir carreras 
facilitando el analisis visual de los resultados, los colores asignados fueron:
\begin{itemize}
	\item MC. Inteligencia Artificial en azul
	\item Ing. Nanotecnologia en amarillo
	\item Ing. Biomedica en verde
	\item Ing. Física en morado
	\item Cruces en rojo.
\end{itemize}

\subsection{Ejecucion del AG}
Con todo lo anterior menc